% ===========================================================================
\chapter{Introduction}
% ===========================================================================

\section{Background}

The past decade has witnessed a fundamental restructuring of
telecommunications infrastructure driven by the convergence of
software-defined networking, network virtualization, and cloud-native
service delivery. Traditional networks were built from vertically
integrated hardware appliances --- firewalls, deep packet inspection
engines, load balancers, WAN optimizers --- each procured, provisioned,
and operated in isolation. This hardware-centric model proved ill-suited
to the scale, heterogeneity, and economic pressures of modern networks.
In response, the community progressively articulated a
\textit{network virtualization} paradigm in which multiple logical
networks coexist on shared physical substrate, decoupling the service
from the infrastructure that realizes it~\cite{chowdhury_survey_2010}.

Network Function Virtualization (NFV) advances this paradigm one step
further by decoupling the \emph{network functions} themselves from
dedicated middleboxes and running them as software instances --- Virtual
Network Functions (VNFs) --- on commodity servers in data centres and at
the network edge. Yi et al.'s comprehensive survey documents how NFV
reduces capital and operational expenditure, improves service agility,
and enables elastic scaling of VNFs in response to traffic
demand~\cite{yi_comprehensive_2018}. The combination of NFV with
software-defined control establishes the foundation on which modern
5G systems are built: the 3GPP 5G system architecture assumes
virtualized, service-based functions orchestrated across a
programmable substrate~\cite{noauthor_5g_nodate}, and this assumption is
deepening as the research agenda moves toward 6G, where global
coverage, artificial intelligence at the network edge, and
end-to-end security become first-class design
goals~\cite{you_towards_2021}.

\section{Service Function Chaining}

In this virtualized substrate, a concrete end-to-end service is
realized by steering traffic through an ordered sequence of VNFs ---
a \textit{Service Function Chain} (SFC). RFC~7498 formulates the
problem statement for service function
chaining~\cite{p_quinn_problem_2015}, and RFC~7665 defines the
reference architecture, introducing the concepts of service function
paths, classifiers, and the Network Service Header used to carry
service metadata along the
chain~\cite{halpern_service_2015}. RFC~8309 situates this work within
the broader IETF framework for service models, clarifying how a
customer-facing service specification is translated into the
network- and device-level configurations that realize
it~\cite{wu_service_2018}.

An SFC request is characterised by an ordered list of VNFs with
per-function resource demands (CPU, memory, storage), together with
aggregate service-level constraints on end-to-end latency, path
bandwidth, required security properties, lifetime, and isolation
requirements. Delivering such a service on the NFV substrate raises
the \textit{SFC placement and embedding} problem: given an incoming
stream of requests, decide --- in real time --- on which substrate
nodes each VNF should be instantiated and along which paths the
inter-function traffic should be routed. The systematic review by
Santos~et~al. shows that SFC placement is consistently NP-hard across
the many scenarios in which it has been studied and that distributed
deployments add further difficulty through unpredictable delays,
resource heterogeneity, and multi-administrative
boundaries~\cite{santos_service_2021}.

\section{The SFC Placement Problem}

Concretely, an SFC placement policy must simultaneously satisfy a
multi-dimensional set of constraints:

\begin{itemize}
  \item \textbf{Resource feasibility.} Each substrate node has finite
  CPU, memory, and storage capacity; the aggregate demand of all VNFs
  placed on a node must not exceed its residual capacity.
  \item \textbf{Latency.} The accumulated propagation, transmission,
  and processing delay along the SFC path must respect the service's
  end-to-end latency budget.
  \item \textbf{Bandwidth.} Every substrate link used between
  consecutive VNFs must have sufficient residual bandwidth for the
  requested traffic rate.
  \item \textbf{Security.} Each substrate node selected for the chain
  must meet (and preferably exceed) the minimum security level
  required by the service.
  \item \textbf{Isolation.} Services that impose hard
  multi-tenancy restrictions must not share substrate nodes with
  other services.
\end{itemize}

Because accepted SFCs hold resources for their lifetime, myopic
decisions quickly fragment the substrate and degrade long-horizon
acceptance. Worse, virtualized SFC deployments are more susceptible to
failure than their hardware counterparts: the taxonomy assembled by
Shahab and Sharma catalogues the many failure modes that now coexist
--- hypervisor crashes, VNF software faults, link congestion, orchestrator
mistakes, and targeted attacks --- all of which must inform the design of a
placement policy that is both efficient and
dependable~\cite{shahab_fault_2025}.

\section{Security Posture as a First-Class Placement Concern}

The very features that make NFV attractive --- shared commodity
hardware, multi-tenancy, rapid reconfiguration, and dynamic orchestration
--- also expand the attack surface of every service deployed on it.
Pattaranantakul, Vorakulpipat, and Takahashi survey the resulting
threat landscape and conclude that SFC security must be treated as a
pervasive, cross-layer concern spanning the NFV Infrastructure, the VNFs
themselves, the SDN control plane, and the orchestration
layer~\cite{pattaranantakul_service_2023}. The pressure is magnified in
5G and 6G networks, where critical services and massive
IoT populations coexist on the same programmable
substrate~\cite{abdelrazek_measuring_2025, you_towards_2021}.

\textit{Security posture}, in the NIST sense, is the status of a
system's information resources and controls together with the
capability to respond to and recover from
threats~\cite{editor_security_nodate}. Assessing posture requires
\textit{security metrics}: Pendleton et al. systematize these into
measures of vulnerability, defense power, threat severity, and
situational context, and argue that a sound security metric must be
both composable across the system and responsive to the attack--defense
interaction~\cite{pendleton_survey_2017}. Abdelrazek et al. specialize
this perspective to 5G/6G, proposing a life-cycle model of network
function security state driven by attack-surface exposure, known
vulnerabilities, and applied controls~\cite{abdelrazek_measuring_2025}.
A placement policy that is ignorant of posture may satisfy every
static constraint while still concentrating risk on overloaded,
recently-compromised, or vulnerability-laden nodes --- with
consequences ranging from SLA violations to cascading service outages.

\section{Limitations of Classical Placement Approaches}

Classical approaches to the SFC placement problem fall into three broad
families. \textit{Exact methods}, typically integer linear programs,
guarantee optimality but scale poorly and are unsuitable for online
operation on large
substrates~\cite{santos_service_2021}. \textit{Heuristic methods} such
as first-fit and best-fit achieve the required throughput but make
locally greedy choices that ignore both inter-request dependencies and
long-horizon resource dynamics. \textit{Meta-heuristic methods} ---
genetic algorithms, particle-swarm optimization, simulated annealing,
ant-colony search --- improve solution quality at the cost of many
evaluations per request, which again limits their use in real-time
orchestration.

A deeper limitation unites all three: none of them \emph{learn} from
the request stream. Each incoming SFC is solved in isolation, and the
operator's hard-won experience of how traffic distributes, where
incidents recur, or which substrate regions tend to become bottlenecks
is not exploited. The survey by Husen, Chaudary, and Ahmad on
requirements of Future Intelligent Networks argues forcefully that
future operations must be \textit{data-driven and adaptive}: machine
learning must be embedded across functions, layers, and administrative
domains if networks are to remain tractable at 5G/6G
scale~\cite{husen_survey_2022}. Moreover, the taxonomy of fault
tolerance in SFCs by Shahab and Sharma shows that resilience is best
achieved by placement strategies that anticipate failure modes ---
something classical, failure-oblivious algorithms
cannot do~\cite{shahab_fault_2025}.

\section{A Deep Reinforcement Learning Approach with Security Posture Awareness}

Deep Reinforcement Learning (DRL) offers a principled answer. An agent
learns a placement policy from interaction with a simulated substrate:
it observes the current network state, chooses a placement action,
receives a reward signal encoding the operator's objectives, and
gradually refines its policy over many episodes. Once trained, the
policy executes in milliseconds and is naturally compatible with the
multi-objective, long-horizon nature of SFC placement.

This thesis develops an \emph{intelligent, security posture aware} DRL
framework for SFC placement on NFV substrates. In contrast to prior DRL
placement work, the framework incorporates security posture into both the constraint mask and the reward signal,
and evaluates placement decisions against a security cost model
rather than an ad-hoc penalty. The agent is a Maskable Proximal Policy
Optimization (Maskable PPO) learner whose action mask enforces every
hard constraint --- resource, latency, bandwidth, security, and
isolation --- at both training and inference time, so that every chain
the agent accepts is feasible by construction. A Graph Neural Network
(GNN) feature extractor endows the policy with topology awareness, so
that the agent can exploit the structural role of each substrate node
in its placement decisions --- a structural awareness that a flat
multi-layer perceptron cannot offer and that prior classical heuristics
cannot exploit.

\section{Research Motivation and Objectives}

The literature reviewed in Chapter~2 reveals a two-sided research
gap. On one side, every existing work that incorporates security into
SFC placement relies on heuristic, meta-heuristic, or
optimization-based methods that solve each request in isolation and
adapt poorly to dynamic traffic patterns and evolving security
conditions. On the other side, every DRL-based SFC placement agent
ignores security requirements entirely --- security plays no role in
the state, reward, or constraint model. The intersection of these two
research directions is empty: no prior work applies DRL to
security-posture-aware SFC placement.

Within this primary gap, the literature exhibits further deficiencies:

\begin{enumerate}
  \item \textbf{Heuristic-only security awareness.} All
  security-aware SFC placement methods are heuristic or
  optimization-based; none learns from the request stream and none
  can adapt placement policy to changing traffic patterns or shifting
  load distributions.
  \item \textbf{Security-blind DRL.} All DRL-based SFC placement
  agents omit security from their state representation, reward
  signal, and constraint enforcement.
  \item \textbf{Soft-penalty constraints.} Existing DRL agents treat
  placement constraints as penalty terms rather than hard guarantees;
  infeasible actions remain reachable during both training and
  inference.
  \item \textbf{Topology-blind policies.} DRL policies typically
  consume flattened feature vectors and fail to exploit the graph
  structure of the substrate.
  \item \textbf{Unfair evaluation.} Comparative evaluations rarely
  equalize the stream of requests seen by different algorithms, so
  apparent performance differences may reflect stream variance rather
  than algorithmic merit.
\end{enumerate}

The objectives of this thesis follow directly from these gaps:

\begin{enumerate}
  \item Bridge the primary gap by designing a \textbf{DRL framework
  for security-posture-aware SFC placement} that learns adaptive
  placement policies from the request stream while treating security
  requirements as first-class constraints.
  \item Develop a \textbf{security cost model}
  that quantifies operational risk as a function of per-node incident
  probability (driven by posture, load, and accumulated incident
  pressure), chain co-location, and time exposure, and use it both
  as a reward signal and as an evaluation metric.
  \item Guarantee \textbf{hard-constraint feasibility} through action
  masking, so that every chain the agent accepts satisfies resource,
  latency, bandwidth, security posture, and isolation requirements by
  construction.
  \item Integrate a \textbf{Graph Neural Network feature extractor}
  into the policy so that placement decisions reflect substrate
  topology and the structural role of each node.
  \item Perform a \textbf{fair comparative evaluation} against a
  risk-minimising best-fit baseline, aligning the random-number
  stream across algorithms so that both see identical request
  sequences.
\end{enumerate}

\section{Thesis Contributions}

The principal contributions of this thesis are:

\begin{enumerate}
  \item A \textbf{Maskable PPO placement agent with hard-masked
  security posture constraints}, in which resource, latency,
  bandwidth, security posture, and isolation requirements are all
  enforced before any action is sampled --- so every chain the agent
  accepts is feasible by construction.
  \item A \textbf{security posture model} capturing per-node incident
  probability driven by effective security score, load, and accumulated
  incident pressure, exposing this risk signal to the GNN policy and
  the Best-Fit baseline.
  \item A \textbf{topology-aware GNN feature extractor} supporting
  GCN, GAT, and GraphSAGE variants with attention-weighted aggregation
  and MLP-based global context fusion, endowing the DRL policy with
  structural awareness of the substrate graph.
  \item A \textbf{windowed acceptance-ratio reward modifier} that
  delivers curriculum-style feedback by scaling terminal rewards
  according to recent acceptance performance, improving training
  stability.
  \item A \textbf{fair comparative evaluation framework} that aligns
  the random-number generator state of the DRL agent and the
  risk-minimising best-fit baseline so that both process identical
  request streams.
  \item A \textbf{comprehensive experimental analysis} covering
  acceptance ratio, security margin, substrate utilization, and
  per-node SFC/VNF tenancy distributions.
\end{enumerate}

\section{Thesis Organization}

The remainder of this thesis is organized as follows.
Chapter~2 surveys the recent literature on SFC placement, covering
security-aware placement, DRL-based placement, and classical heuristic
and optimization-based approaches, and situates the proposed framework
within this landscape. Chapter~3 formalizes the system model,
introduces the Markov Decision Process formulation, and defines the
security posture risk model. Chapter~4 presents the proposed
methodology in detail, including the Maskable PPO agent, the GNN
feature extractor, the reward design, the action-masking
implementation, and the best-fit baseline. Chapter~5 describes the
experimental setup and presents comparative results, and Chapter~6
concludes the thesis and outlines directions for future work.
